\documentclass{beamer}
\usepackage[citestyle=authoryear, style=authoryear, backend=biber, uniquename=init, giveninits=true]{biblatex}
\usepackage{lmodern}
\usepackage{amsmath}
\usepackage{mathtools}
\usepackage[version=3]{mhchem} % Formula subscripts using \ce{}, e.g., \ce{H2SO4}
\usepackage{siunitx}
\usepackage[utf8]{inputenc}
% https://tex.stackexchange.com/a/155611/56227
\usepackage{textcomp}
% footnote per page https://tex.stackexchange.com/a/1088/56227
\usepackage{perpage}
\usepackage{hyperref}

\graphicspath{{./Figures/}}

% clear title and doi from citations
% from https://tex.stackexchange.com/q/165481/56227
\AtEveryCitekey{\clearfield{doi}\clearfield{title}}

% use dash for range phrase by default, and one unit
\sisetup{range-phrase=--, range-units = single}

% turn off nav bar
\beamertemplatenavigationsymbolsempty

% footnote per page https://tex.stackexchange.com/a/1088/56227
\MakePerPage{footnote}

% symbols for footnotes
\renewcommand{\thefootnote}{\fnsymbol{footnote}}

%smaller footnotes, adapted from http://tex.stackexchange.com/a/146021
\setbeamerfont{footnote}{size=\fontsize{6pt}{0pt}}

% numeric footnote symbols
\def\thefootnote{\arabic{footnote}}

% set themes
\usetheme{Frankfurt}
\usecolortheme{seahorse}
\usefonttheme{serif}

% adjust footnote width
%\newlength{\tmpwidth}
%\setlength{\tmpwidth}{0.65\paperwidth}
%\addtobeamertemplate{footnote}{\hsize\tmpwidth}{}

%footnote spacing avoids navigation / margins, adapted from http://tex.stackexchange.com/a/44231
%\addtobeamertemplate{footnote}{\vspace{-6pt}\advance\hsize-0.5cm}{\vspace{6pt}}
%\makeatletter
% Alternative A: footnote rule
%\renewcommand*{\footnoterule}{\kern -3pt \hrule \@width 2in \kern 8.6pt}
% Alternative B: no footnote rule
% \renewcommand*{\footnoterule}{\kern 6pt}
%\makeatother

%make frame for each section
\AtBeginSection[]{
  \begin{frame}
  \vfill
  \centering
  \begin{beamercolorbox}[sep=8pt,center,shadow=true,rounded=true]{title}
    \usebeamerfont{title}\insertsectionhead\par%
  \end{beamercolorbox}
  \vfill
  \end{frame}
}

%auto incrementing titles from http://tex.stackexchange.com/a/231533
\newcounter{expensive}
\newcommand\ExpTitle{%
  \frametitle{\refstepcounter{expensive}{Chemical kinetic integration is \textbf{expensive}~--}~\theexpensive}}
\resetcounteronoverlays{expensive}

% define background template env
\newenvironment{background}{%
\usebackgroundtemplate{%
\rule{0pt}{\paperheight}%
\hspace*{\paperwidth}%
\makebox[0pt][r]{\includegraphics[width=35mm]{logo}}%
}}{}

% space between paragraphs beamer
% https://tex.stackexchange.com/a/22644/56227
\def\parend/{\\~\\}

\bibliography{presentation.bib}

%opening
\title{Accelerating reacting flow simulations via vectorized chemical kinetic integration}
\author{Nick Curtis}
\institute{University of Connecticut}
\date{\today}

\begin{document}

\begin{background}
\maketitle
\end{background}

\begin{background}
\section{Introduction}
\end{background}

\begin{frame}
\frametitle{Improving combustion technology}
The push to design new clean, efficient combustion technologies has driven both:
\begin{itemize}
 \item the development of combustion devices operating in new regimes such as low-temperature combustion as well as,
 \item the creation of large detailed chemical kinetic models for a wide range of fuel classes from bioalcohols~\footfullcite{SARATHY2009852} to gasoline surrogates\footfullcite{SARATHY201867}.
\end{itemize}
These next generation combustion devices are often controlled by chemical processes, rather than current methods which rely on directly physical processes.\parend/
Fuel selection and composition are key parameters for the design of such technologies, and the use of realistic chemical modeling is critical in order to achieve predictive numerical simulations.
\end{frame}

\begin{frame}
 \frametitle{Challenges of reactive-flow simulations}
 \textbf{Model size:} Chemical kinetic models for transportation and energy relevant fuels may consist of hundreds to thousands of chemical species, with potentially tens of thousands of reactions (e.g., gasoline \footfullcite{Mehl:2011cn} and jet fuel\footfullcite{Naik2011434})\parend/
 \textbf{Range of scales:} A typical reactive-flow simulation encompasses a wide range of physical and temporal scales, e.g., the flow-through time of a combustion chamber is typically on the order of seconds to milliseconds, while chemical timescales may be just picoseconds in duration\footfullcite{Lu:2009gh} due to the presence of highly reactive radicals and associated short chemical timescales.\parend/
\end{frame}

\begin{frame}
 \frametitle{Numerical stiffness}
 The large range of time-scales present in these chemical kinetic models---known as numerical stiffness---has traditionally been handled by use of implicit integration techniques, which require repeated evaluation and factorization of the chemical kinetic Jacobian.
 \begin{itemize}
  \item Naive implementations of these operations scale \textbf{quadratically} and \textbf{cubically} with the number of species in a model, respectively\footfullcite{Lu:2009gh}.
 \end{itemize}
 However, using even modestly sized chemical kinetic models with such a solver can incur severe computation cost (e.g.,~\footfullcite{Moiz2016123}) for realistic reactive-flow simulations.
\end{frame}


\begin{frame}
 \frametitle{Reducing the cost of chemical kinetics}
  A host of techniques have been developed to reduce the cost of chemical kinetic calculations while maintaining fidelity~\footfullcite{turanyi2016analysis}, e.g.:
  \begin{itemize}
   \item removal of unimportant species and reactions,
   \item lumping of species with similar thermochemical properties, and
   \item time-scale methods that reduce numerical stiffness
  \end{itemize}
  Today we will focus on two such (complementary) techniques, analytical Jacobian evaluation and vectorized chemical kinetic integration.
\end{frame}

\begin{background}
\section{Analytical jacobian evaluation}
\end{background}

\begin{frame}
 \footfullcite{HANSEN2018257}
\end{frame}


\begin{background}
 \section{Vectorized ODE integration}

\begin{frame}
 \frametitle{Operator splitting and vectorization}
 Most reacting-flow codes utilize the operator-splitting techniques (e.g.~\footfullcite{Knio:1999}$^{,}$~\footfullcite{Ren:2008}), separating a large system of coupled partial differential equations (PDEs), such that different physical processes are solved independently.
 \begin{itemize}
  \item Independent initial value problem (IVP) of chemical kinetic ordinary differential equations (ODEs) at each computational location (e.g. cell) in the domain~$\rightarrow$~vectorization!
 \end{itemize}
\end{frame}

\end{background}


\end{document}